\documentclass[11pt]{article}
\usepackage[T1]{fontenc}
\usepackage{lmodern,amsmath,amssymb,amsthm,mathtools,microtype}
\usepackage[margin=1.05in]{geometry}
\usepackage[hidelinks]{hyperref}
\hypersetup{pdftitle={A mod-three obstruction to repeated XOR of small-bias sources}}
\newtheorem{theorem}{Theorem}
\newtheorem{lemma}[theorem]{Lemma}
\newtheorem{proposition}[theorem]{Proposition}
\newtheorem{corollary}[theorem]{Corollary}
\newcommand{\F}{\mathbb F_3}
\newcommand{\E}{\mathbb E}
\newcommand{\wt}{\operatorname{wt}}
\title{A mod-three obstruction to repeated XOR of small-bias sources}
\author{}
\date{}
\begin{document}
\maketitle
\vspace{-3em}
\begin{abstract}
For $q=3^m$ and fixed $t$, every XOR of up to $t$ points from our
construction has Hamming weight divisible by three. The source has
bias at most $\exp(-q^{c_t/(1+\log^*q)})$.
At any prescribed bias $q^{-C}$ there is a polynomial-time sampler
using $O_{t,C}(\log q)$ unbiased bits. With a longer seed the support
condition persists for $\Omega(\sqrt{\log q})$ copies at the same
inverse-polynomial bias. The distinguisher is the width-three
permutation program that counts ones modulo three.
\end{abstract}

\section{Results and construction}
For a bit string $Y$ indexed by $x\in\F^m$, write $S_Y(x)=(-1)^{Y_x}$.
A distribution on $Y$ has bias at most $\delta$ if
$|\mathbb E\prod_{x\in T}S_Y(x)|\le\delta$ for every nonempty set
of indices $T$. The signs denote real numbers when taking expectations
and elements of $\F$ when evaluating polynomials. We use
$\log^*$ for the iterated logarithm with base three.

\begin{theorem}\label{thm:main}
Let $t\ge1$ be a fixed integer and $q=3^m$. For all sufficiently
large $m$ there are distributions on $q$ bits for which
\begin{equation}
 \wt(y_1\oplus\cdots\oplus y_j)\equiv0\pmod 3
 \qquad(1\le j\le t,\; y_1,\dots,y_j\in\operatorname{supp}(Y)).
 \label{eq:checksum}
\end{equation}
One can obtain either of the following guarantees.
\begin{enumerate}
\item\label{item:tiny} The bias is at most
$\exp(-q^{c_t/(1+\log^*q)})$ for some $c_t>0$.
The sampler uses polynomial time and polynomially many unbiased bits.
\item\label{item:short} For each fixed $C>0$ the bias is at most $q^{-C}$.
The sampler runs in polynomial time and uses $O_{t,C}(\log q)$ unbiased bits.
\end{enumerate}
At bias $q^{-C}$ there is also a polynomial-time, polynomial-randomness
sampler satisfying \eqref{eq:checksum} with
$t=\lfloor c\sqrt m\rfloor$ for an absolute $c>0$ and all sufficiently
large $m$ depending on $C$.
\end{theorem}

A three-state program adds each input bit to its current state modulo
three, starts at zero, and accepts states one and two. It rejects every
word in \eqref{eq:checksum}. At $q=3^m$, $m\ge1$, its uniform acceptance
probability is $2/3+2^{1-q}/3$. Indeed, with
$\zeta=e^{2\pi i/3}$ the uniform probability of residue zero is
$[1+2\operatorname{Re}(((1+\zeta)/2)^q)]/3$.
Thus the support of the XOR even misses a language of density greater
than $2/3$. All lengths can be covered by taking the largest power of
three $q\le n$, appending independent small-bias bits, and letting the
program ignore that suffix. The pure Hamming-weight conclusion in
\eqref{eq:checksum} applies to the original $q$ coordinates.

Lee and Viola~\cite[Theorem 1.4(iv)]{LV} obtained a mod-three linear
distinguisher for $t$ copies at bias $q^{-\Omega(1/t)}$. Here the test
counts ordinary Hamming weight for every $t$, and the exponent $C$ in
part~\ref{item:short} is arbitrary. The bound in part~\ref{item:tiny} is
weaker than $\exp(-\Omega(q))$, the bias in the two-copy self-dual-code
example. These counterexamples concern assertions uniform over
all distributions of the given bias.

\paragraph{The support condition.}
Suppose each sign word has a polynomial representative of total degree
at most $d$ in the $m$ index trits. Whenever $td<2m$ we have
\[
  \sum_{x\in\F^m}\prod_{i=1}^j S_{y_i}(x)=0\quad\text{in }\F
               \qquad (1\le j\le t).
\]
Indeed, each monomial of degree less than $2m$ has an exponent zero or
one in some coordinate. Summation in that coordinate gives zero in
$\F$. The left side is $q-2\wt(y_1\oplus\dots\oplus y_j)$ modulo
three. This proves \eqref{eq:checksum}, including for repeated support
points. We will multiply two independent random sign words. Their
ternary degrees add, while their real parity biases multiply.

\paragraph{Interpolation on few points.}
For $1\le K\le q$ let $d_K=2\lceil\log_3K\rceil$. A uniformly random
ternary polynomial on $\F^m$, sampled by choosing independent coefficients
on all monomials of degree at most $d_K$ and individual degree at most
two, has independent uniform values on any specified $K$ points.
Indeed for each point separately a random linear map to
$\F^{\lceil\log_3 K\rceil}$ separates it from the other points with
positive probability by the union bound. The indicator of its fiber is
a polynomial of degree at most $d_K$, by multiplying factors
$1-(\ell(x)-a)^2$. These indicators interpolate all possible
assignments, so the linear evaluation map is surjective.
\section{Nearly balanced signs of low ternary degree}
\label{sec:multi-coin}

Throughout this section, a sign polynomial is a polynomial over
$\mathbb F_3$ taking only the values $1,-1$. Its mean denotes the real
mean of the corresponding signs on uniform, independent ternary inputs.
Let $\log^* L$ be the number of iterated base-three logarithms needed to
reach a value at most one.
\begin{lemma}\label{lem:multi-coin}
For every integer $L\ge1$, one can construct a sign polynomial $P$
in $O(L)$ ternary variables with real mean $-3^{-a}$ for some integer
$L\le a\le 4L$. Its degree is $O(1+\log^* L)$.
Its description and evaluation take time polynomial in $L$.
There is also a direct construction of a sign polynomial of degree
$O(\log(2L))$ and absolute mean at most $3^{-L}$ on $O(L)$ variables.
\end{lemma}

\begin{proof}
For $j\ge1$, take $Q_j=\sum_{i=1}^j u_i v_i$ with all pairs
independently uniform over $\mathbb F_3^2$. For $c\in\mathbb F_3$ put
$H_{j,c}=2\mathbf1[Q_j=c]-1$. This is a sign polynomial because
$\mathbf1[z=c]=1-(z-c)^2$. Conditional on the $u_i$ not all being
zero, $Q_j$ is uniform. Otherwise $Q_j=0$. Consequently
\[
    \mathbb E[H_{j,0}]-\mathbb E[H_{j,1}]=2\cdot 3^{-j}.
\]
For the direct construction, let $R=3^h=2T+1$ be the least power
of three at least $2L+1$. Give one of the $R$ possible addresses
in $\mathbb F_3^h$ the constant value $-1$. Give the remaining addresses
the values $H_{j,0}$ and $-H_{j,1}$ for $1\le j\le T$,
using the same list of $T$ independent pairs for all the branches.
The resulting mean is
\[
    R^{-1}\left(-1+\sum_{j=1}^T2\cdot3^{-j}\right)
       =-3^{-T}/R .
\]
An address indicator has degree at most $2h$, so the degree is at most
$2h+4$.

For the improved degree, begin with a sign of mean $-1/3$, obtained by
returning $1$ on one of the three values of a single ternary variable.
We describe a step starting with a sign of mean $-3^{-a}$. Choose any
integer $0\le h\le a$ and set $R=3^h$. Use an independent uniform
address $s\in\mathbb F_3^h$. Regard its coordinates as ternary digits
$[s_i]\in\{0,1,2\}$ and put
\[
   J=a-h+1+\sum_{i=1}^h 3^{i-1}[s_i].
\]
We can realize a single polynomial $Q$ of degree at most four whose
conditional distribution at each address is that of $Q_J$.
Use $a-h+1$ unconditional summands $u v$. In addition, for digit $i$
use two disjoint groups of $3^{i-1}$ fresh pairs. Multiply the sum of
the products in one group by $\mathbf1[s_i\ne0]$, and the sum in the
other by $\mathbf1[s_i=2]$. These two indicators have degree at most
two over $\mathbb F_3$. At any address the number of active pairs is
exactly $J$.

Use a fresh ternary selector. On its three possible values, return
$2\mathbf1[Q=0]-1$, $-(2\mathbf1[Q=1]-1)$, or the old sign,
respectively. Since $J$ is uniform on the integers from $a-h+1$
through $a-h+R$, the mean is
\[
 \frac13\left(
  \frac{2}{R}\sum_{j=a-h+1}^{a-h+R}3^{-j}-3^{-a}\right)
 =-3^{-(a+R+1)}.
\]
Thus the exponent advances from $a$ to $a+3^h+1$. If the old degree
was $d$, the new degree is at most $\max\{10,d+2\}$. The number of
additional variables and arithmetic operations is $O(a+R)$.

For completeness, here is a choice of steps for every target $L$.
For $1\le L\le 16$ one may apply the step with $h=0$ until its
exponent first reaches $L$. For $L>16$ first recurse with target
$h=\lceil\log_3 L\rceil$, obtaining an exponent $a\in[h,4h]$.
Apply the step with this $h$. The resulting exponent lies between
$L$ and $4L$ since $3^h<3L$ and $4h+1\le L$ for $L>16$.
There are $O(1+\log^* L)$ recursive stages. The sum of the variable
counts over the stages is $O(L)$, since each successive recursive
target is at most half the preceding target whenever that target
exceeds $16$. The same construction evaluates the polynomial by
arithmetic circuits with $O(L)$ gates up to an absolute constant.
\end{proof}
\section{A short-seed source for parities on large sets}
\label{sec:large-parities}

Use $q=3^m$ coordinates indexed by $x\in\mathbb F_3^m$.
Fix $2\le g\le m$ and partition the $m$ trits into $g$ chunks
of sizes at most $\lceil m/g\rceil$. Write $x_a$ for the value of
chunk $a$. We construct random signs $H(x)$ with the following
two properties. Every output is a polynomial over $\mathbb F_3$
of total degree at most $8\lceil m/g\rceil$. For every nonempty
$T\subseteq\mathbb F_3^m$ and $0<\delta<1$,
\begin{equation}
 \left|\mathbb E\prod_{x\in T}H(x)\right|
 \le 2\delta+\exp\!\bigl(-2^{-g}|T|^{1/g}\bigr).
 \label{eq:large-parity-bound}
\end{equation}
The explicit sampler uses $O_g(\log q+\log(1/\delta))$ unbiased bits.

For each $i\in[g]$ construct a table of bits $b^{(i)}_y$ indexed by
values of chunk $i$. Use a $\delta$-biased source for this table.
Independently construct a second bit string containing $u^{(i)}_y$
and all lookup entries $L^{(i)}_{y,a}(v)$, where $a\ne i$
and $v$ is a value of chunk $a$. Use an explicit generator for
degree-$g$ polynomials over $\mathbb F_2$ whose error in the
expectation of each sign $(-1)^P$ is at most $\delta$.
Viola's theorem gives this generator by XORing $g$ independent
sufficiently small-bias sources~\cite{Viola}.
Take the choices for distinct $i$ independently. Regarding the
table entries as elements of $\mathbb F_3$, set
\[
 P_i(x)= b^{(i)}_{x_i}+u^{(i)}_{x_i}
       +\sum_{a\ne i}L^{(i)}_{x_i,a}(x_a),\qquad
 P(x)=\sum_{i=1}^{g} P_i(x),\qquad
 H(x)=(-1)^{\mathbf 1[P(x)=0]}.
\]
Any lookup on two chunks is a ternary polynomial of total degree at
most $4\lceil m/g\rceil$, by interpolation with individual degree at
most two. In $\mathbb F_3$ the sign $H$ equals $-1-P^2$.
This gives the degree bound for every seed. There are at most
$O(g^2q)$ table entries. The seed bound follows from the small-bias
and degree-$g$ generators.

To prove \eqref{eq:large-parity-bound}, choose $i$ so that projection
onto chunk $i$ gives $M\ge |T|^{1/g}$ distinct values on $T$. Such
a chunk exists because $T$ lies in the Cartesian product of its
projections. Fix the tables belonging to all other chunks. Their
contribution to $P$ is now a fixed function $P_0(x)$. Suppress
the superscript $(i)$ from the remaining tables and let
$T_y=\{x\in T:x_i=y\}$. Work in $\mathbb F_4$ with a primitive cube
root of unity $\omega$ and write $\operatorname{Tr}$ for the trace to
$\mathbb F_2$. For $z\in\mathbb F_3$ we have the binary identity
$\mathbf 1[z=0]=1+\operatorname{Tr}(\omega^z)$. Define
\[
 V_y=\sum_{x\in T_y}\omega^{P_0(x)}
                   \prod_{a\ne i}\omega^{L_{y,a}(x_a)},\qquad
 c_y=\operatorname{Tr}\!\left((1+\omega)\omega^{u_y}V_y\right).
\]
Conditional on $u,L$, the parity of the zero count of $P$ on $T$ is
an affine binary function of $b$ with coefficients $c_y$. Indeed,
$\omega^{b_y}=1+(1+\omega)b_y$ for binary $b_y$. Thus the desired
sign expectation has absolute value at most
$\delta+\Pr_{u,L}[c_y=0\text{ for all occupied }y]$.

Under independent uniform $u,L$, the $c_y$ for distinct fibers are
independent. For each occupied fiber, treat the lookup factors as
formal variables $z_{a,v}$. The polynomial $V_y$ has distinct
monomials for distinct $x\in T_y$, each with a nonzero coefficient
and one variable from each of $g-1$ groups. With independent
uniform substitutions from $\{1,\omega\}$ it is nonzero with
probability at least $2^{-(g-1)}$. In detail, fix any one monomial
and expose every variable outside it. Its coefficient remains
nonzero because no monomial strictly contains it. The resulting
multiaffine polynomial is nonzero on at least one of the
$2^{g-1}$ points of the remaining grid, by univariate interpolation
in each variable. If $V_y\ne0$, at least one of the two values
$u_y=0,1$ makes $c_y=1$, since $1,\omega$ span $\mathbb F_4$ over
$\mathbb F_2$ and the trace pairing is nondegenerate. It follows
that the probability all $c_y$ vanish is at most
$(1-2^{-g})^M\le \exp(-2^{-g} M)$ under uniform $u,L$.

Finally, each $c_y$ is a binary polynomial of degree at most $g$
in $u,L$. Each factor $\omega^v$ is affine over $\mathbb F_4$ in
a binary lookup entry $v$, and the trace is binary linear. Hence
every binary linear combination of the $c_y$ has degree at most
$g$. The indicator that all $M$ coefficients vanish is the average
of $(-1)^{\sum_y a_y c_y}$ over $a\in\mathbb F_2^M$. The generator
for $u,L$ therefore changes its expectation by at most $\delta$.
This bound is uniform in the fixed $P_0$ and proves
\eqref{eq:large-parity-bound}.

\subsection{Hashing to reach parities of logarithmic size}
\label{sec:hashed-large-parities}

One can replace the coordinate projection above by a ternary linear
hash without increasing the asymptotic degree. This gives an
$O_{g,C}(\log q)$-bit sampler for all sets of size at least
$O_{g,C}(\log q)$ at error $q^{-C}$.

Fix $s=\lfloor m/g\rfloor$. Identify the index grid with
$\mathbb F_{3^m}$ in a fixed ternary basis, and define
$h_\lambda(x)$ to be the first $s$ coordinates of $\lambda x$.
Draw $R$ independent hash parameters. In the summand for hash $h$
use $y=h(x)$ and the lookup entries $L_{y,a}(x_a)$ for
\emph{all} $g$ of the fixed chunks $a$, including any chunks
intersecting the domain of the hash. Thus the summand is
\[
 b_{h(x)}+u_{h(x)}+\sum_{a=1}^g L_{h(x),a}(x_a).
\]
Use independent table samplers for the $R$ summands, with the
degree-$(g+1)$ generator for the $u,L$ tables. All $s$ hash
coordinates are ternary linear forms in the index. A table entry
$L_{h(x),a}(x_a)$ can be interpolated in the $s$ hash coordinates
and at most $\lceil m/g\rceil$ other coordinates, with individual
degree two before substitution. The final output sign has degree
at most $4(s+\lceil m/g\rceil)\le8\lceil m/g\rceil$.

For any $K\le |T|$, fix $K$ distinct indices of $T$. For a
uniform $\lambda$, two distinct such indices have equal hashes
with probability $3^{-s}$, since multiplication by their
nonzero difference permutes $\mathbb F_{3^m}$. The probability
that none of the $R$ hashes separates all $K$ indices is at most
$\bigl(\binom K2 3^{-s}\bigr)^R$. Conditional on any good list
of hashes, choose the first separating hash and condition on the
tables of all other summands. At least $K$ of its fibers are
occupied by $T$. The fiber argument above applies to all points
of $T$, including those outside the selected $K$. Indeed, the
full tuple $(x_1,\ldots,x_g)$ determines the index, so the
lookup monomials for distinct points in a fiber are distinct even
for a non-coordinate hash. Now each monomial has $g$ factors.
The same interpolation argument gives
$\Pr[c_y\ne0]\ge2^{-(g+1)}$ for an occupied fiber under
uniform $u,L$. Each $c_y$ has binary degree at most $g+1$.
Writing $\eta$ for the error of each table sampler in the
corresponding sign tests, we obtain
\[
 \left|\mathbb E\prod_{x\in T}H(x)\right|
 \le 2\eta+\exp(-2^{-(g+1)}K)+
      \left(\binom K2 3^{-s}\right)^R
\]
for ideal uniform hashes.

For an unbiased bit seed, approximate a uniform field index by
reducing a uniform $b$-bit integer modulo $q$, with statistical error at
most $q2^{-b}$. Using $b=\lceil\log_2(8Rq/\delta)\rceil$ for
each hash changes the expectation of any sign by at most
$\delta/4$. At $\delta=q^{-C}$ set $\eta=\delta/8$,
$K=\lceil 2^{g+1}\ln(4/\delta)\rceil$, and choose a fixed
integer $R>2g(C+1)$. For sufficiently large $q$, these choices
give bias at most $\delta$ for every $|T|\ge K$, since
$\binom K2 3^{-s}\le q^{-1/(2g)}$. All $R$ hash parameters and
table seeds together cost $O_{g,C}(\log q)$ bits. A field model
and its multiplication can be computed in time polynomial in $q$,
for instance by enumerating and testing monic degree-$m$
polynomials over $\mathbb F_3$ until an irreducible one is found.
\section{A short-seed source for parities on small sets}
\label{sec:small-parities}

Let $q=3^m$, let $1\le K\le q$, and let $0<\delta\le1/2$.
We index coordinates by $x\in\mathbb F_3^m$. Put
$d=2\lceil\log_3 K\rceil$ and $v=\lceil\log_3(4/\delta)\rceil$.
The sampler below satisfies the following three bounds.
\begin{enumerate}
\item For every seed its output is a sign polynomial in $x$ of degree
at most $d(2\lceil\log_3(2v+1)\rceil+4)$.
\item For every nonempty $S$ of at most $K$ locations,
$|\mathbb E\prod_{x\in S}L(x)|\le\delta$.
\item It uses $O(\log(2q)+K+\log(1/\delta))$ unbiased random bits.
It runs in time polynomial in $q$ and $1/\delta$ when
$K=O(\log(1/\delta))$.
\end{enumerate}

Use the $D\le q$ reduced monomials of degree at most $d$.
By the interpolation argument in Section~1, a uniformly
random coefficient vector gives independent uniform trits at any fixed
$k\le K$ distinct locations.

Set $h=\lceil\log_3(2v+1)\rceil$, $R=3^h=2T+1$. Use one ternary
polynomial for each of $h$ address digits, each of $T$ variables
$u_1,\ldots,u_T$, and each of $T$ variables $v_1,\ldots,v_T$.
All use the monomials above. Sample the three concatenated coefficient
vectors for addresses, $u$'s, and $v$'s independently, each from a
ternary $\eta$-biased source, with $\eta$ chosen below. At location
$x$, interpret the address digits as an integer $a\in[0,2T]$.
Return $-1$ if $a=T$. If $a=T-j$, return $H_{j,0}$, and if $a=T+j$,
return $-H_{j,1}$, where
\[
 Q_j=\sum_{b=1}^j u_b(x)v_b(x),\qquad
 H_{j,c}=2\mathbf1[Q_j=c]-1.
\]
The direct construction of Lemma~\ref{lem:multi-coin} gives the
stated degree upon composition. Under uniformly random coefficients,
all gadget variables on $S$ are independent, and each sign has mean
$-3^{-T}/R$. Every nonempty such parity therefore has absolute mean
at most $3^{-v}\le\delta/4$.

We record a fact for comparing the true and pseudorandom coefficients.
An affine subspace of a ternary vector space has indicator equal to an
average of additive characters times unit complex phases. Thus its
probability changes by at most $\eta$ when uniform inputs are replaced
by $\eta$-biased inputs, meaning that every nontrivial additive
character has absolute expectation at most $\eta$. A decision tree
making at most $B$ adaptive ternary linear queries has at most $3^B$
leaves. Each leaf is an affine subspace or empty, so the expectation of
an output in $[-1,1]$ changes by at most $3^B\eta$. If the original
tree may run for more than $B$ queries with uniform probability $p$,
truncate it at $B$, returning zero on an unfinished branch. Its timeout
probability under the biased source is at most $p+3^B\eta$, by applying
the same argument to the timeout indicator. Altogether, the change in
expectation for the original tree is at most $2p+2\cdot3^B\eta$.

Fix a test on $k=|S|\le K$ locations. First condition on the address
and $u$ pools. The output product is a bounded function of at most $k$
linear forms in the $v$ coefficient pool, one $Q_j$ per nonconstant
branch. The depth-$k$ tree which reads these forms shows that replacing
this pool by uniform coefficients changes the test by at most $3^k\eta$.

Next fix the addresses and integrate over this uniform $v$ pool.
At distinct sites the selected forms use independent arrays of
uniform trits. Each selected $Q_j$ is uniform unless the first $j$
$u$ entries at that site vanish, in which case $Q_j=0$. The conditional
mean of the product is therefore a product of bounded factors. Each
factor can be computed by reading the first $j$ entries of its $u$
list, stopping at the first nonzero. Across the $k$ sites under uniform
$u$ coefficients, the query counts are independently dominated by a
geometric variable $G$ with $\Pr[G>r]=3^{-r}$ for integers $r\ge0$.
In particular $\mathbb E[2^G]=4$. Thus the total number exceeds $B$
with probability at most $4^k2^{-B}$. Every query is a linear form in
the $u$ coefficients. The decision-tree bound permits the replacement
of this pool by uniform coefficients.

Finally, keep the $v$ pool integrated out and condition on the now
uniform $u$ pool. Let $J$ be the first nonzero position at a given
site, with $J=T+1$ if the list vanishes. Its factor depends only on
whether $a<T$, $a=T$, or $a>T$, and whether $|a-T|\ge J$.
This function has at most four changes between consecutive addresses.
Read the ternary digits of $a$ from most significant to least until
the current prefix interval has constant factor. After $r$ digits,
the chance that this still needs another digit under a uniform address
is at most $4\cdot3^{-r}$, since each cut can intersect the interior
of at most one prefix interval. If $C$ is the query count, the tail
sum identity gives $\mathbb E[2^C]\le1+4\sum_{r\ge0}(2/3)^r=13$.
The addresses on $S$ are independent under uniform coefficients. The
total address query count consequently exceeds $B$ with probability
at most $13^k2^{-B}$, uniformly in the conditioned $u$ pool. All
address digit queries are linear forms in the address coefficients.
Applying the decision-tree bound one last time replaces the address
pool by uniform coefficients.

Take $B=\lceil K\log_2 13+\log_2(32/\delta)\rceil$ and
$\eta=\delta/(32\cdot3^B)$. The three hybrid errors together with the
uniform bias sum to at most $\delta$. Each pool has $M\le O(Dv)$
ternary coefficients. For an explicit sampler, choose an extension
$\mathbb F_{3^s}$ of size at least $4M/\eta$, with $s$ minimal.
For independent uniform $b,z$ in this field take the $M$ trits
$\operatorname{Tr}_{\mathbb F_{3^s}/\mathbb F_3}(b z^i)$,
$0\le i<M$. A nonzero character has bias at most $(M-1)/3^s$ by
averaging over $b$ and counting roots of a nonzero polynomial in $z$.
Using $\lceil\log_2(8\cdot3^{2s}/\eta)\rceil$ unbiased bits and
reduction modulo $3^{2s}$ to choose $(b,z)$ changes a character
expectation by at most $\eta/4$. The resulting source is $\eta$-biased
and uses $O(\log(2M)+\log(1/\eta))$ bits. The field and its arithmetic
can be constructed by exhaustive search in time polynomial in $3^s$.
This proves the seed and time bounds.
\section{Putting the bounds together}
\label{sec:assembly}
Fix $g=100t$. For the first part of Theorem~\ref{thm:main}, let $L$
be a parameter and put
\[
   K=\left\lceil(2^g(2L+2))^g\right\rceil .
\]
Use Lemma~\ref{lem:multi-coin} to take a sign $F$ of degree
$O(1+\log^*L)$ on $v=O(L)$ trits with absolute real mean at most
$3^{-L}$. Independently choose $v$ uniform ternary polynomials
$A_1,\ldots,A_v$ of degree at most $d_K$. Set
$L_0(x)=F(A_1(x),\ldots,A_v(x))$. On any nonempty $T$ of size
at most $K$, its signs are independent with common mean of magnitude
at most $3^{-L}$. Its ternary degree satisfies
\[
       \deg L_0\le d_K\deg F
          =O_t\big(\log(2L)(1+\log^* L)\big).
\]
Take an independent large-set sign $H$ from
Section~\ref{sec:large-parities}, using truly independent unbiased
table entries. The proof of \eqref{eq:large-parity-bound} gives
$|\mathbb E\prod_{x\in T}H(x)|\le\exp(-2^{-g}|T|^{1/g})$.
For $|T|>K$ this is at most $e^{-2L}$. Hence $L_0H$ has bias
at most $\max(3^{-L},e^{-2L})$.

Take $L=3^{\lfloor m/(A_t(1+\log^*m))\rfloor}$ with
$A_t$ a sufficiently large integer. Then $K\le q$ for all sufficiently
large $m$. Also $\log^* L\le 1+\log^*m+O(1)$, and the constant in
the low-degree bound is independent of $m$. Taking $A_t$ large
enough yields $t\deg L_0\le m$. Since $t\deg H\le
8t\lceil m/g\rceil<m/2$ for large $m$, the product satisfies the
support condition. There are at most $q$ coefficients per $A_i$,
so the number of uniform trits needed is polynomial in $q$.
They can all be sampled to within total variation
$e^{-2L}/2$ by reducing a sufficiently long uniform binary integer
modulo the cardinality of their product space. This adds at most
$e^{-2L}$ to any parity bias, and every possible seed still gives
the degree bound. Brute-force monomial enumeration and evaluation
take polynomial time in $q$. Decreasing a positive constant $c_t$
absorbs the rounding terms and proves part~\ref{item:tiny}.

For part~\ref{item:short}, set $\delta=q^{-C}$ and
$K=\lceil 2^{g+1}\ln(4/\delta)\rceil$.
Multiply the independent signs from Sections~\ref{sec:small-parities}
and~\ref{sec:hashed-large-parities}. The first has bias at most
$\delta$ for $1\le |T|\le K$, and the second has bias at most
$\delta$ for $|T|\ge K$. Their seed lengths sum to
$O_{t,C}(\log q)$. Their total ternary degree is at most
\[
     O_{t,C}((\log m)^2)+8\lceil m/g\rceil < 2m/t
\]
for sufficiently large $m$. This proves the short-seed assertion.
\section{A growing number of summands}
\label{sec:growing-summands}

Fix $C>0$. For all sufficiently large $m$ and every integer
$1\le t\le \sqrt m/250$, there is a $q^{-C}$-biased distribution
on $q=3^m$ bits such that the XOR of any $1$ through $t$ support
points has Hamming weight divisible by three. The sampler runs in
polynomial time in $q$ for fixed $C$. This construction uses a
polynomial number of random bits. The logarithmic seed bound in the
fixed-$t$ construction does not enter this argument.

Put $g=100t$, $b=\lfloor m/g\rfloor$, $a=\lceil m/g\rceil$,
$K=\lceil(C+2)m\rceil$, and $M=\lfloor3^{b/3}\rfloor$.
In particular, $K<M$ for all sufficiently large $m$, uniformly in
the stated range of $t$. Index coordinates by $x\in\mathbb F_3^m$.
Use the independent product of the following two random sign functions.
For the first sign $L(x)$, use the small-set construction of
Section~\ref{sec:small-parities} with cutoff $K$ and error
$q^{-C}/4$. Its degree is $O_C(\log^2 m)$ for every seed.

For the second sign, take a uniformly random reduced ternary
polynomial $P$ of total degree at most $D=2(a+b)$. More precisely,
the coefficients of all monomials of degree at most $D$ and
individual degrees at most two are independent uniform trits.
Set $H(x)=(-1)^{\mathbf1[P(x)=0]}$, represented over $\mathbb F_3$
by $-1-P(x)^2$. For every $K\le |T|\le M$, the evaluations of $P$
on $T$ are independent uniform trits. Indeed, each point in $T$ is
separated from the others by a linear map to
$\mathbb F_3^{\lceil\log_3 |T|\rceil}$, and the indicator of the
fiber containing this point has degree at most
$2\lceil\log_3 |T|\rceil\le D$, as in
Section~1. Thus
\[
  \left|\mathbb E\prod_{x\in T}H(x)\right|
   =3^{-|T|}\le q^{-C-2}.
\]

For $|T|>M$, choose $M$ distinct points in $T$. Some linear map
$h:\mathbb F_3^m\to\mathbb F_3^b$ is injective on these points,
since a random linear map collides on a pair with probability at
most $\binom M2 3^{-b}<1$. This map can be chosen separately in
the analysis of each $T$. Partition the $m$ coordinates into $g$
chunks of size at most $a$, and use the correction polynomial
\[
   W(x)= B_{h(x)}+U_{h(x)}+\sum_{j=1}^g A_{h(x),j}(x_j),
\]
with independent uniform binary table entries embedded in
$\mathbb F_3$. Here $x_j$ denotes the entire $j$th chunk. Since
each table entry depends on at most $a+b$ ternary linear forms,
$W$ has reduced degree at most $D$. For $P_0$ an independent
uniform polynomial from the same space as $P$, we have
$P_0+W\overset{\mathrm{law}}=P$. Conditional on $P_0$, the
uniform-table calculation in Section~\ref{sec:hashed-large-parities}
applies. At least $M$ fibers are occupied, the table entries in
different fibers are independent, and each occupied fiber has
probability at least $2^{-(g+1)}$ of making its parity unbiased
over the unused binary entry $B_y$. Hence
\[
   \left|\mathbb E\prod_{x\in T} H(x)\right|
       \le \exp\bigl(-2^{-(g+1)}M\bigr).
\]
This is at most $q^{-C}/4$ for all sufficiently large $m$ uniformly
in $t$. Indeed, $g\le(2/5)\sqrt m$ and
$b\ge(5/2)\sqrt m-1$, so
\[
 \log\bigl(2^{-(g+1)}M\bigr)
 \ge \left(\tfrac56\log 3-\tfrac25\log 2\right)\sqrt m-O(1).
\]
There are at most $q$ coefficients to sample in $P$. Reducing a
uniform binary integer modulo $3^{\#\mathrm{coefficients}}$
with $\lceil (\#\mathrm{coefficients})\log_2 3+
\log_2(16q^C)\rceil$ random bits incurs total variation distance at
most $q^{-C}/16$. Thus all the parity bounds suffice with a fixed
polynomial-length unbiased binary seed. Independence of $L$ and $H$
makes their biases multiply on each test set. For sets smaller than
$K$ use the bound on $L$, and for all larger sets use the bound on
$H$.

Each sign $LH$ has degree at most $D_*=O_C(\log^2 m)+4(a+b)$.
We have $tD_*\le O_C(\sqrt m\log^2m)+8m/100+4t<2m$ for large
$m$. Every product of at most $t$ such signs is a polynomial of
degree less than $2m$. Its sum over the ternary grid vanishes in
$\mathbb F_3$, because each monomial has some coordinate of
exponent zero or one, and the sum in that coordinate is zero.
Writing the sign product as $(-1)^z$ yields
$0=\sum_x(-1)^{z(x)}=q-2|z|$ in $\mathbb F_3$. This is the
claimed checksum for the XOR. Counting weight modulo three is a
standard-order width-three permutation program.

\subsection{An upper bound for an exact mod-three checksum}
Suppose a $\delta$-biased distribution on $q$ bits has support $A$, and
every XOR of exactly $t$ support words, repetitions allowed, has the same
Hamming weight modulo three. For $0\le\delta<1$,
\[
               t\le \frac{\log q}{\log(2/(1+\delta))}.
\]
In particular $t\le(1+o(1))\log_2 q$ when $\delta=o(1)$.

Translate $A$ and its distribution by $v\in A$. The resulting
support $A'$ contains zero and has the same bias bound. Set $u=v$
if $t$ is odd, and $u=0$ if $t$ is even. For any XOR $z$ of at most
$t$ words in $A'$, padding with zero puts $z\oplus u$ in the original
$t$-fold sumset. Taking $z=0$ identifies its common residue as $|u|$.
With $w_i=(-1)^{u_i}$ we obtain
$\sum_i w_i z_i= |z\oplus u|-|u|=0\pmod 3$.

Sample $t$ independent words from $A'$. Write $p_i\in\{-1,1\}^t$
for their signs at coordinate $i$ and $\mathbf{1}$ for the all-one
pattern. For every $J\subseteq[t]$, the preceding weighted identity
for the XOR on $J$ gives
$\sum_i w_i\prod_{j\in J}p_{ij}=\sum_i w_i$ in $\mathbb F_3$.
By linearity, for every multilinear ternary polynomial $Q$,
\[
       \sum_i w_i Q(p_i)=\Big(\sum_i w_i\Big)Q(\mathbf{1})
                     \qquad\text{in }\mathbb F_3.
\]
Fix one coordinate $i$. For $k\ne i$, two-coordinate bias gives
$\Pr[p_k=p_i]\le((1+\delta)/2)^t$. Singleton bias gives the same
bound on $\Pr[p_i=\mathbf{1}]$. If $q((1+\delta)/2)^t<1$, some
choice of words has $p_i$ distinct from all other patterns and from
$\mathbf{1}$. The indicator of this chosen pattern is the multilinear
polynomial $Q(X)=\prod_{j=1}^t(1+p_{ij}X_j)/2^t$ over $\mathbb F_3$.
Its left side above is $w_i\ne0$, and its right side is zero.
This contradiction proves the bound.

\begin{thebibliography}{9}
\bibitem{LV} Chin Ho Lee and Emanuele Viola.
\emph{Some limitations of the sum of small-bias distributions}.
Theory of Computing 13(16), 1--23, 2017.
\href{https://doi.org/10.4086/toc.2017.v013a016}{doi:10.4086/toc.2017.v013a016}.
\bibitem{Viola} Emanuele Viola.
\emph{The sum of $D$ small-bias generators fools polynomials of degree $D$}.
Computational Complexity 18, 209--217, 2009.
\href{https://eccc.weizmann.ac.il/report/2007/132/}{ECCC TR07-132}.
\end{thebibliography}
\end{document}
